\textbf{\textit{Soal. }}In an isosceles triangle $ABC$, $AB=AC$. Take a point $O$ inside the triangle (not on the boundary). Draw a circle $\omega$ with center $O$ and passing through $C$. Suppose $\omega\cap BC=\{C,D\}$, $\omega\cap AC=\{C,E\}$. Denote the circumcircle of $\triangle AEO$ as $\Gamma$. And suppose $\Gamma\cap\omega=\{E,F\}$. Prove that the circumcenter of $\triangle BDF$ is on $\Gamma$.

\begin{proof}[\textbf{Solusi. } (Monday, 10 October 2022) ]
    Misalkan $K$ adalah perpotongan kedua antara $ED$ dengan $\Gamma$. Akan dibuktikan bahwa $K$ adalah pusat dari lingkaran luar segitiga $BDF$.

Thanks EulersTurban for the asymptote picture :)
\begin{center}
    \begin{asy}
    /* Geogebra to Asymptote conversion, documentation at artofproblemsolving.com/Wiki go to User:Azjps/geogebra */
    import graph; size(12cm); 
    real labelscalefactor = 0.8; /* changes label-to-point distance */
    pen dps = linewidth(0.7) + fontsize(10); defaultpen(dps); /* default pen style */ 
    pen dotstyle = black; /* point style */ 
    real xmin = -16.581447502554504, xmax = 14.190483526605371, ymin = -6.724424998317599, ymax = 11.674458045226984;  /* image dimensions */
    
     /* draw figures */
    draw((0,8.08)--(-3.54,0), linewidth(0.4) + blue); 
    draw((-3.54,0)--(3.54,0), linewidth(0.4) + blue); 
    draw((3.54,0)--(0,8.08), linewidth(0.4) + blue); 
    draw(circle((-4.009081052926384,4.068417484239327), 5.671465874858991), linewidth(0.4) + linetype("4 4")); 
    draw((1.6578143075369038,4.296062258503339)--(-0.9271496309748504,-0.6925933202094363), linewidth(0.4) + blue); 
    draw((-0.9271496309748504,-0.6925933202094363)--(-1.4869827036317487,0), linewidth(0.4) + blue); 
    draw((-3.54,0)--(-0.9271496309748504,-0.6925933202094363), linewidth(0.4) + blue); 
    draw((-0.9271496309748504,-0.6925933202094363)--(-2.5134913518158757,-1.4022987957026576), linewidth(0.4) + blue); 
    draw((-2.5134913518158757,-1.4022987957026576)--(1.6578143075369038,4.296062258503339), linewidth(0.4) + blue); 
    draw((0,8.08)--(-0.9271496309748504,-0.6925933202094363), linewidth(0.4) + blue); 
    draw((-0.9271496309748504,-0.6925933202094363)--(3.54,0), linewidth(0.4) + blue); 
    draw((-3.54,0)--(-2.5134913518158757,-1.4022987957026576), linewidth(0.4) + blue); 
    draw((0,8.08)--(-2.5134913518158757,-1.4022987957026576), linewidth(0.4) + blue); 
    draw((0,8.08)--(1.026508648184126,1.459133764181308), linewidth(0.4) + blue); 
    draw((-2.5134913518158757,-1.4022987957026576)--(1.026508648184126,1.459133764181308), linewidth(0.4) + blue); 
    draw((-1.4869827036317487,0)--(1.026508648184126,1.459133764181308), linewidth(0.4) + blue); 
    draw((1.026508648184126,1.459133764181308)--(-0.9271496309748504,-0.6925933202094363), linewidth(0.4) + blue); 
    draw((1.026508648184126,1.459133764181308)--(1.6578143075369038,4.296062258503339), linewidth(0.4) + blue); 
    draw(circle((1.0265086481841261,1.4591337641813078), 2.9063224386545805), linewidth(0.4) + linetype("4 4") ); 
    draw((1.026508648184126,1.459133764181308)--(3.54,0), linewidth(0.4) + blue); 
     /* dots and labels */
    dot((0,8.08),dotstyle); 
    label("$A$", (0.0900294910503877,8.279904295926729), NE * labelscalefactor); 
    dot((-3.54,0),dotstyle); 
    label("$B$", (-3.465213193128005,0.20528531830126717), NW * labelscalefactor); 
    dot((3.54,0),linewidth(4pt) + dotstyle); 
    label("$C$", (3.6251860583690156,0.1651130845817375), E * labelscalefactor); 
    dot((1.026508648184126,1.459133764181308),dotstyle); 
    label("$O$", (1.1144214508983992,1.651485732204335), NE * labelscalefactor); 
    dot((-1.4869827036317487,0),linewidth(4pt) + dotstyle); 
    label("$D$", (-1.416429273431982,0.1651130845817375), NW * labelscalefactor); 
    dot((1.6578143075369038,4.296062258503339),linewidth(4pt) + dotstyle); 
    label("$E$", (1.737091073551112,4.463542092571411), NE * labelscalefactor); 
    dot((-0.9271496309748504,-0.6925933202094363),linewidth(4pt) + dotstyle); 
    label("$F$", (-0.854018001358564,-0.5379010055100315), S * 3 * labelscalefactor); 
    dot((-2.5134913518158757,-1.4022987957026576),linewidth(4pt) + dotstyle); 
    label("$K$", (-2.4408212332799937,-1.2409150956018005), S * 3* labelscalefactor); 
    clip((xmin,ymin)--(xmin,ymax)--(xmax,ymax)--(xmax,ymin)--cycle); 
     /* end of picture */
\end{asy}
\end{center}

Perhatikan bahwa $OF=OD=OE=OC$ yang merupakan jari-jari lingkaran $\omega$. Karena $A,E,O,F$ di lingkaran $\Gamma$ dan $OE=OC$ maka $\dangle ACO = \dangle OEC = \dangle OEA = \dangle OFA$. Lalu, karena $OC=OF$ maka $\dangle OCF = \dangle CFO$. Hal ini menyebabkan $\dangle ACF = \dangle ACO + \dangle OCF = \dangle OFA + \dangle CFO = \dangle CFA$ sehingga didapat $AC=AF$ atau lebih jauh $AB=AC=AF$. Fakta ini juga otomatis menyebabkan $\triangle OCA \cong \triangle OFA \implies \dangle OAC = \dangle FAO$.

Sekarang, perhatikan karena $O$ pusat $\omega$, maka $\dangle DOF = 2\dangle DEF$. Karena $E,A,F,K$ berada di satu lingkaran $\Gamma$, maka $\dangle DEF = \dangle KEF = \dangle KAF = \dangle KOF$. Ini berarti $\dangle DOF = 2\dangle KOF \implies \dangle DOK = \dangle KOF$. Namun, kita tahu bahwa $DO=OF$, maka didapatkan bahwa $KO \perp DF$ sehingga didapatkan $DOFK$ adalah belah ketupat dengan $DK=DF$ dan $\dangle FKO = \dangle OKD$.

Karena $AC=AB=AF$, maka $A$ adalah pusat lingkaran luar segitiga $BFC$, sehingga kita punya $\dangle BAC = 2\dangle BFC$ dan $\dangle BAF = 2\dangle BCF$. Tapi karena $D,E,F,C$ berada di lingkaran $\omega$, maka $\dangle BCF = \dangle DCF = \dangle DEF = \dangle KOF$. Dari sini karena $K,A,E,F$ berada di lingkaran $\Gamma$ juga, akan didapat $\dangle BAF = 2\dangle KOF = 2\dangle KAF$ yang menyebabkan $\dangle BAK = \dangle KAF$. Namun, karena $AB=AF$, maka $BAFK$ adalah belah ketupat dengan $BK=KF$. Dari sini kita mendapat $BK=KF=KD$, yang berarti $K$ haruslah menjadi pusat lingkaran luar segitiga $BDF$. Terbukti.



\end{proof}